\section{Related Work}
\label{sec:related_work}

Our research intersects five major areas of empirical software engineering and software supply-chain security: structural ecosystem measurement, malicious package detection, continuous-integration workflow analysis, registry governance and provenance, and temporal dependency reconstruction.

\subsection{Ecosystem Structure and Centrality-Based Prioritization}
Pioneering empirical studies established that open-source package registries exhibit extreme structural asymmetry and small-world properties. Zimmermann et al.~\cite{zimmermann2019} analyzed npm's dependency topology, revealing that a tiny core of central packages and maintainers transitively affects a vast majority of the ecosystem. Zahan et al.~\cite{zahan2022} study metadata weak-link signals, while Liu et al.~\cite{liu2022} examine dependency-tree vulnerability propagation.

Zimmermann et al.~\cite{zimmermann2019} already studied maintainer-union reach, greedy coverage, and vetting. Wang~\cite{wang2026} evaluates package-node protection in cascade simulations. Our contribution is an empirical comparison of coverage, account-origin exposure, and all-source exclusion at workflow-associated repository endpoints. We evaluate established baselines and a stated heuristic, without claiming a new coverage algorithm or demonstrating that PageRank fails in an experiment where it was not evaluated.

\subsection{Malicious Package Benchmarks and Detection Limits}
A parallel line of security research focuses on detecting malicious packages published via typosquatting, dependency confusion, and account takeovers~\cite{ohm2020,sejfia2022}. Benchmark evaluations, such as Guo et al.'s extensive study of npm malware detectors~\cite{guo2026}, show that static analysis tools face severe limitations due to obfuscation, dynamic runtime execution, and environment-keyed payloads. 

Because reactive post-publication malware detection suffers from non-trivial evasion rates, proactive hardening of legitimate, high-leverage maintainer accounts is a complementary research direction~\cite{schorlemmer2024,okafor2025}. Our work complements malware detection by optimizing how defenders allocate proactive security interventions before credentials or tokens are compromised.

\subsection{Continuous-Integration Workflows and Reachability Limits}
Software supply-chain risks manifest concretely in downstream CI/CD pipelines. Koishybayev et al.~\cite{koishybayev2022} and Muralee et al.~\cite{muralee2023} demonstrated that GitHub Actions workflows frequently execute third-party build scripts with access to repository secrets, cloud deployment credentials, and package publish tokens. Recent incidents motivate artifact-level checks: Railway's preliminary report describes potential secret exposure affecting at most 1.85\% of workspaces, approximately 100,000, and notes that builds had resumed~\cite{railway2026}. LiteLLM reported a release workflow redirected to attacker-controlled source, with separate approval and log-deletion credentials~\cite{litellm2026}. These are contextual examples rather than observations in our npm graph.

Crucially, recent empirical studies have emphasized that package-level dependency relationships do not equal actual code execution. Weeraddana, Alfadel, and McIntosh~\cite{weeraddana2024} studied CI time spent updating unused npm dependencies. Shen et al.~\cite{gao2025emse} reported that 99\% of client vulnerabilities in their studied corpus arose from dependencies and 81.26\% of package-level findings were false positives under finer analysis. These corpus-specific results do not establish registry-wide false-positive rates. Jayasuriya et al.~\cite{jayasuriya2024} observed similar granularity effects when analyzing syntactic breaking changes across versions. In our work, we address this fundamental reachability gap by evaluating both coarse package-name projections and version-continuous graphs, treating topological reach as a conditional review surface rather than installed-code evidence.

\subsection{Registry Governance, Cryptographic Provenance, and Signing}
In response to account compromises, registries and standards bodies have introduced cryptographic signing and provenance attestations. Schorlemmer et al.~\cite{schorlemmer2024} audited signing across Maven, PyPI, Docker Hub, and Hugging Face, finding that adoption remains minimal and unevenly distributed. Kalu et al.~\cite{okafor2025} interviewed industry practitioners, revealing significant operational friction in managing cryptographic keys and configuring automated pipeline attestations. Bloom Security's SlugFest disclosure concerns mutable repository-name bindings before a first trusted publish, rather than a break of OIDC cryptography~\cite{npmtrustedpub2026}. GitHub announced expiry of unvalidated bindings on October 2~\cite{npmexpiry2026}; reported exposure is not evidence that all affected packages were maliciously published. Pillar's October 19 OpenCode report describes a September 28 dependency fix without a CVE~\cite{opencode2026}, illustrating a coverage gap for CVE-indexed detection.

Simultaneously, He, Vasilescu, and K{\"a}stner~\cite{he2025} demonstrated that local dependency version pinning is insufficient to protect clients against upstream attacks, as adversaries exploit unpinned sub-dependencies and registry-level re-resolutions. These findings motivate examining account controls alongside version and artifact evidence; they do not establish control efficacy in our graph.

\subsection{Temporal Reconstruction and Empirical Standards}
A critical methodological vulnerability in supply-chain studies is future-leakage bias---evaluating historical attacks using contemporary registry graphs. Pinckney et al.~\cite{pinckney2023} introduced \texttt{npm-follower} to mirror registry event streams, while Rahman et al.~\cite{rahman2026} released \texttt{CHRONO-RESOLUTION} to reconstruct exact dependency trees at release points. Shafin and Ahmed~\cite{shafin2026} highlighted the necessity of time-indexed forecasting.

Our reporting is informed by the ACM SIGSOFT Empirical Standards~\cite{ralph2020}. By avoiding pseudo-replication---specifically rejecting the treat\-ment of fixed-graph nested budgets or repeated random orders as independent sampling units---our study provides reproducible descriptive comparisons and exploratory scalarization results.

\begin{table*}[t]
\centering
\small
\caption{Selected literature and the scope of our empirical investigation relative to major related literature.}
\label{tab:lit_positioning}
\begin{tabularx}{\textwidth}{lYYYY}
\toprule
\textbf{Research Dimension} & \textbf{Structural Topology \& Centrality} & \textbf{Malicious Package Detectors} & \textbf{CI/CD \& Function Call-Graphs} & \textbf{This Study (Q1 Journal Expansion)} \\
\midrule
\textbf{Focus Literature} & Zimmermann~\cite{zimmermann2019}, Zahan~\cite{zahan2022}, Wang~\cite{wang2026} & Ohm~\cite{ohm2020}, Sejfia~\cite{sejfia2022}, Guo~\cite{guo2026} & Koishybayev~\cite{koishybayev2022}, Weeraddana~\cite{weeraddana2024}, Shen~\cite{gao2025emse} & Structural objectives ($C, H, P, J_\alpha$), version continuity, field availability \\
\midrule
\textbf{Prioritization Metric} & Dependency and maintainer reach; coverage; metadata signals & Reactive post-event incident signatures & Taint analysis on workflows, call-graph reachability & Submodular coverage, modular exposure reduction, supermodular all-source exclusion \\
\midrule
\textbf{Graph Representation} & Dependency graphs and resolved trees & Isolated package tarballs & Dynamic execution traces, static ASTs & Dual evaluation: Package-projected vs. Version-continuous closures (42.4\% pair diff) \\
\midrule
\textbf{Maintainer Risk Prior} & Uniform equal weight & Unweighted & Unweighted & Heterogeneous weights: Solo bottlenecks ($H_{\text{sole}}$), log-portfolio ($H_{\text{log}}$), package-level sole reach ($H_{\text{pkg-sole}}$) \\
\midrule
\textbf{Temporal Grounding} & Static single snapshot & Retrospective positive-only datasets & Static GitHub Actions scripts & 45 root records (44 distinct tuples), 1,848 bounded edge attempts ($t_{\text{cutoff}} < t_{\text{incident}}$), 90-day prospective cohort \\
\bottomrule
\end{tabularx}
\end{table*}
